\documentclass[10pt,a4paper]{amsart}
\usepackage[margin=19mm]{geometry}
\usepackage{amsmath,amssymb,mathtools}
\usepackage[scr=rsfs,cal=euler]{mathalfa}
\usepackage{xfrac}
\usepackage[unicode,hidelinks,bookmarks=false]{hyperref}
\newcommand{\ie}{i.e.\ }
\newcommand{\cf}{cf.\ }
\newcommand{\resp}{respectively\ }
\newcommand{\Subsection}{\subsection}
\providecommand{\nobreakdash}{\nobreak\hbox{-}}
\setlength{\emergencystretch}{2em}
\multlinegap0pt
\usepackage{bm}
\usepackage{bbm}
\usepackage{dsfont}
\usepackage{xspace}
\usepackage{cancel}
\usepackage{mathtools}
\usepackage{amsthm}
\makeatletter
\@ifundefined{lemma}{\newtheorem{lemma}{Lemma}}{}
\@ifundefined{theorem}{\newtheorem{theorem}[lemma]{Theorem}}{}
\makeatother
\makeatletter
\newcommand{\Cost}{\mathrm{Cost}}
\expandafter\ifx\csname claim\endcsname\relax
\newenvironment{claim}[1]{\par\noindent\underline{Claim:}\space#1}{}
\fi
\expandafter\ifx\csname claimproof\endcsname\relax
\newenvironment{claimproof}[1]{\par\noindent\underline{Proof:}\space#1}{\hfill $\blacksquare$}
\fi
\newcommand{\fix}{\mathrm{fix}}
\newcommand{\mc}{\mathcal}
\newcommand{\wt}{\widetilde}
\makeatother
\title[Transportation cost and contraction coefficient for channels]{Transportation cost and contraction coefficient for channels on von Neumann algebras}
\author{Roy Araiza, Marius Junge, Peixue Wu}
\date{Source: arXiv:2506.04197v1 (CC BY 4.0)}
\begin{document}
\maketitle

\noindent\emph{Source and licence.} Transportation cost and contraction coefficient for channels on von Neumann algebras. Version arXiv:2506.04197v1 (CC BY 4.0), distributed under the Creative Commons Attribution 4.0 International licence, as recorded in the arXiv licence metadata for this exact version and shown on the abstract page https://arxiv.org/abs/2506.04197v1. Full rights notice: licenses/arxiv-2506.04197-CC-BY-4.0.txt.\par\smallskip

\noindent\textit{Editorial scope.} Counted unit: the Theorem whose source label is \texttt{thm:expected\_length} in the article (source0.tex lines 1780--1785, source label \texttt{thm:expected\_length}), delivered together with the article's own complete original proof of it (source0.tex lines 1787--1852, one \texttt{proof} environment of 66 lines). No mathematics is written, repaired or re-derived: statement and proof are the article's own text line for line. Required context retained uncounted: def:von Neumann algebra finite group (lines 1676--1678); def:Lipschitz norm finite group (lines 1698--1700); lemma:calculation finite group (lines 1703--1715). Every definition, display and label of those lines is retained unchanged. Omitted from this package and not redistributed: 1--1675, 1679--1697, 1701--1702, 1716--1779, 1786--1786, 1853--2623. Nothing inside a retained excerpt is omitted or altered: every retained body is delivered exactly as the article's lines are written, so no omission receipt of any kind accompanies this package. Citations: the retained excerpts carry no citation command at all, so no bibliography entry of the article is transcribed and no reference is redistributed. Reference adaptation: none was needed and none was made; every reference inside the delivered unit resolves inside it, so no unresolved reference or citation is produced. Printed slips kept literal: no printed word, symbol, hypothesis or number of the counted statement or of its proof was altered. Layout and class: the article's own class is \texttt{amsart} with packages including geometry, amsmath, amsxtra, amssymb, amsthm, amsfonts, mathtools, mathrsfs, multirow, tikz, quantikz, physics, bbm, subfigure; the wrapper is the lane's pinned \texttt{amsart} editorial wrapper at 10pt on A4, which restates the theorem-like environments the retained text needs and adds the article's own macro definitions that the retained text needs (\texttt{\textbackslash Cost}, \texttt{\textbackslash fix}, \texttt{\textbackslash mc}, \texttt{\textbackslash wt}), with extra wrapper packages (bm, bbm, dsfont, xspace, cancel, mathtools). The delivered numbering is the unit's own. Assets: no retained span contains a figure environment, a graphics-inclusion command, a PDF-page inclusion, a table or a TikZ picture; no image file is copied into this package, the package declares no asset dependency and no image file is redistributed with it. Licence: version arXiv:2506.04197v1 of this article is distributed under the Creative Commons Attribution 4.0 International licence, as recorded in the arXiv licence metadata for this exact version and shown on the abstract page https://arxiv.org/abs/2506.04197v1. A full rights notice is delivered at licenses/arxiv-2506.04197-CC-BY-4.0.txt. Characters: every retained excerpt is pure ASCII, so the delivered engine, XeLaTeX with its default Latin Modern text font, reports no missing character.
\begin{equation}\label{def:von Neumann algebra finite group}
    \mathcal{N} = L_{\infty}(G,\mu) := \{ f : G \to \mathbb{C} \mid f \text{ bounded} \}, \quad \mathcal{H} = L_2(G,\mu),
\end{equation}

\begin{equation}\label{def:Lipschitz norm finite group}
    ||| f |||_{\mathcal{S}} := \sup_{s \in S} \|[\lambda_s, \pi(f)]\|_{op}.
\end{equation}

\begin{lemma}\label{lemma:calculation finite group}
    Let $G$ be a finite group with the normalized counting measure $\mu$, and let $\mc H = L_2(G,\mu)$. Suppose $f \in \mc N = L_{\infty}(G,\mu)$, where $(\pi,\mc H)$ is the standard representation defined by $\pi(f)(\wt f) = f\wt f$ for $\wt f \in \mc H$. Then:
    \begin{enumerate}
        \item We have
        \[
            \|\pi(f)\|_{op} = \|f\|_{\infty} = \sup_{g \in G} |f(g)|.
        \]
        \item For any symmetric generating set $S \subseteq G$ (i.e.\ $S=S^{-1}$), we have
        \[
            |||f|||_{\mc S} = \sup_{s\in S} \sup_{g \in G}|f(sg) - f(g)|.
        \]
    \end{enumerate}
\end{lemma}

\begin{theorem}\label{thm:expected_length}
    Let $\mc N$, $\mc H$, and $\mc S$ be as in the finite group setup \eqref{def:von Neumann algebra finite group}, with the Lipschitz semi-norm \eqref{def:Lipschitz norm finite group}. Then we have
    \[
        \Cost_{\mc S}(E_{\fix})= \int_G \ell_{S}(g)\, d\mu(g).
    \]
\end{theorem}

\begin{proof}
\textbf{Lower bound:} Denote $\mc N_{\mathrm{s.a.}}$ as the set of real-valued functions. We need to show that 
\[
    \sup_{f \in \mc N_{\mathrm{s.a.}}, |||f|||_{\mc S} \le 1} \|\pi(E_{\fix}(f) - f)\|_{op} \ge \int_G \ell_{S}(g) \, d\mu(g).
\]

Consider the function 
\[
    f_{\mc S} := \ell_{S} - \left(\int_G \ell_{S}(g)\, d\mu(g)\right) 1_G.
\]
Note that $f_{\mc S}$ is real-valued and hence self-adjoint. By construction, $E_{\fix}(f_{\mc S}) = 0$. We now estimate its Lipschitz semi-norm:
\[
    |||f_{\mc S}|||_{\mc S} = \sup_{s \in S} \sup_{g \in G} |f_{\mc S}(sg) - f_{\mc S}(g)| = \sup_{s \in S} \sup_{g \in G} |\ell_{S}(sg) - \ell_{S}(g)|.
\]
Since adding or removing a single generator changes the word length by at most $1$, we have 
\[
    |\ell_{S}(sg) - \ell_{S}(g)| \le 1.
\]
Thus $|||f_{\mc S}|||_{\mc S} \le 1$.

Using Lemma~\ref{lemma:calculation finite group}, we then have
\[
    \|\pi(E_{\fix}(f_{\mc S}) - f_{\mc S})\|_{op} = \|\pi(-f_{\mc S})\|_{op} = \|f_{\mc S}\|_{\infty} = \sup_{g \in G} |f_{\mc S}(g)|.
\]
Since $f_{\mc S}(e) = \ell_{S}(e) - \int_G \ell_{S}(g)\, d\mu(g) = 0 - \int_G \ell_{S}(g)\, d\mu(g)$, we see that
\[
    \|f_{\mc S}\|_{\infty} \ge |f_{\mc S}(e)| = \int_G \ell_{S}(g)\, d\mu(g).
\]
This establishes the desired lower bound.

\medskip

\textbf{Upper bound:} Now we show
\[
    \sup_{f \in \mc N_{\mathrm{s.a.}}, |||f|||_{\mc S} \le 1} \|\pi(E_{\fix}(f) - f)\|_{op} \le \int_G \ell_{S}(g) \, d\mu(g).
\]

Take any $f \in \mc N_{\mathrm{s.a.}}$ with $|||f|||_{\mc S} \le 1$. Using the definition of $E_{\fix}$ and the invariance of $\mu$, we have:
\begin{align*}
    \|\pi(E_{\fix}(f) - f)\|_{op} &= \sup_{g \in G} \left| f(g) - \int_G f(h) \, d\mu(h) \right| \\
    &= \sup_{g \in G} \left| f(g) - \int_G f(hg) \, d\mu(h) \right| \\
    &\le \sup_{g \in G} \int_G |f(g) - f(hg)| \, d\mu(h).
\end{align*}

For a fixed $h \in G$, write $h = s_1 s_2 \cdots s_{\ell_S(h)}$ with $s_i \in S$. A telescoping sum argument yields
\[
    |f(g) - f(hg)| \le \sum_{k=1}^{\ell_S(h)} |f(u_k) - f(s_k u_k)|,
\]
where $u_k := s_{k+1}\cdots s_{\ell_S(h)}g$. Since $|||f|||_{\mc S} \le 1$, we have
\[
    |f(u_k) - f(s_k u_k)| \le \sup_{u \in G} |f(u) - f(s_k u)| \le |||f|||_{\mc S} \le 1.
\]
Thus
\[
    |f(g) - f(hg)| \le \ell_{S}(h).
\]
Substituting this back, we get
\[
    \|\pi(E_{\fix}(f) - f)\|_{op} \le \int_G \ell_{S}(h) \, d\mu(h).
\]
Since this holds for all $f$ with $|||f|||_{\mc S} \le 1$, we conclude
\[
    \sup_{f \in \mc N_{\mathrm{s.a.}}, |||f|||_{\mc S} \le 1} \|E_{\fix}(f) - f\|_{op} \le \int_G \ell_{S}(g) \, d\mu(g).
\]
Combining the lower and upper bounds completes the proof.
\end{proof}

\end{document}
